\documentclass[11pt,a4paper]{article}
\usepackage[margin=2.4cm]{geometry}
\usepackage{amsmath,amssymb,booktabs,hyperref,microtype}
\usepackage[T1]{fontenc}
\usepackage{lmodern}
\hypersetup{colorlinks,linkcolor=black,urlcolor=blue!50!black}
\setlength{\parskip}{4pt}
\newcommand{\dd}{\,\mathrm{d}}
\newcommand{\Hs}{H_s}

\title{\texttt{ww3recon}: reconstructing $E(f,\theta)$ from WW3 partition parameters\\[4pt]
\large Technical note}
\author{2026 wave coupling project}
\date{2026-09-16, package version 0.2}

\begin{document}
\maketitle

\section{Problem}

WAVEWATCH III archives such as IOWAGA \texttt{GLOBMULTI\_ERA5\_GLOBCUR\_01} store, at every grid
cell and time, up to six spectral partitions $k=0,\dots,5$ (0 = wind sea, 1--5 = swell systems
ordered by energy) with five bulk parameters each, plus whole-spectrum integrals:
\begin{center}
\begin{tabular}{llll}
\toprule
archive field & symbol & meaning & used for\\
\midrule
\texttt{phs}$_k$ & $\Hs^{(k)}$ & partition significant height [m] & partition energy $m_0^{(k)}$\\
\texttt{ptp}$_k$ & $T_p^{(k)}$ & partition peak period [s] & peak frequency $f_p^{(k)} = 1/T_p^{(k)}$\\
\texttt{pdir}$_k$ & $\bar\theta^{(k)}$ & partition mean direction (\emph{coming-from}) & directional centre\\
\texttt{pspr}$_k$ & $\sigma_\theta^{(k)}$ & partition directional spread [deg] & directional width\\
\texttt{pws}$_k$ & $w^{(k)}$ & wind-sea fraction in the partition & family blend\\
\texttt{hs}, \texttt{t0m1}, \texttt{t02} & $\Hs$, $T_{m0,-1}$, $T_{m02}$ & whole-spectrum integrals & closure check, shape fit, tail fit\\
\bottomrule
\end{tabular}
\end{center}
Storage is about 1/30 of the full spectrum. The task is to rebuild a directional variance density
$E(f,\theta)$ on the native grid (36 frequencies, $f_i = 0.0339\,\mathrm{Hz}\times1.1^{\,i}$; 24 directions at
$15^\circ$) that reproduces the archived integrals and is as close as possible to the archived 2D spectrum.

\section{Grid, moments and conventions}

Frequency bin widths follow WW3's own rule for a geometric ladder with ratio $r$,
\begin{equation}
\Delta f_i = \tfrac12 f_i\,(r - 1/r), \qquad \Delta\theta = 2\pi/24 .
\end{equation}
Spectral moments and the bulk parameters are discrete sums on this grid,
\begin{equation}
m_n = \sum_{i,j} f_i^{\,n}\,E(f_i,\theta_j)\,\Delta f_i\,\Delta\theta,\qquad
\Hs = 4\sqrt{m_0},\quad T_{m0,-1} = \frac{m_{-1}}{m_0},\quad T_{m02} = \sqrt{\frac{m_0}{m_2}} .
\end{equation}
All archived integrals are truncated at the grid edge ($0.95\,$Hz), so every moment in the fit is
evaluated on the same truncated grid; comparing a reconstruction integrated on a finer or extended
grid to an archived value is not a fair comparison.

\paragraph{Direction conventions.} The archive's \texttt{pdir}$_k$, \texttt{dir} and \texttt{wnddir}
are \emph{coming-from} directions. The 2D spectrum axis is \emph{going-to}
(\texttt{sea\_surface\_wave\_to\_direction}, nautical degrees clockwise from north, stored in
descending order). Partition directions are rotated $\bar\theta^{(k)} \mapsto \bar\theta^{(k)} + 180^\circ$
before use, and the spectrum's direction axis is sorted ascending. Both are confirmed from file
metadata, not inferred.

\section{Model}

Each partition is a separable product of a frequency shape and a directional distribution, and the
partitions are summed:
\begin{equation}
E(f,\theta) = \sum_{k} m_0^{(k)}\, S^{(k)}(f)\, D^{(k)}(\theta),
\qquad m_0^{(k)} = \frac{\big(\Hs^{(k)}\big)^2}{16},
\label{eq:model}
\end{equation}
with $S^{(k)}$ normalised so that $\sum_i S^{(k)}(f_i)\Delta f_i = 1$ and $D^{(k)}$ so that
$\sum_j D^{(k)}(\theta_j)\Delta\theta = 1$. Consequently
$m_0 = \sum_k m_0^{(k)}$ exactly, i.e.\ $\Hs$ is conserved partition by partition by construction,
and the sum is invariant to the partition labelling, so index swapping between neighbouring cells
does no harm. Absent partitions (NaN in the archive) are skipped.

\subsection{Frequency shape}

The shape of partition $k$ is a blend of a wind-sea family $S_{\mathrm{ws}}$ and a swell family
$S_{\mathrm{sw}}$, both peaked at $f_p^{(k)}$ and normalised to unit $m_0$, weighted by the wind-sea
fraction:
\begin{equation}
S^{(k)} = \omega^{(k)} S_{\mathrm{ws}}(f;\,f_p^{(k)},\,\Omega_c) + \big(1-\omega^{(k)}\big) S_{\mathrm{sw}}(f;\,f_p^{(k)},\,\lambda),
\qquad
\omega^{(k)} = \operatorname{clip}\!\Big(\frac{w^{(k)}-0.3}{0.7-0.3},\,0,\,1\Big).
\end{equation}
Blending the normalised shapes rather than switching on a threshold keeps time series continuous when
\texttt{pws} crosses the threshold. Two shape parameters, $\Omega_c$ and $\lambda$, are shared by all
partitions of a point and are the only free quantities; they are determined in Section~\ref{sec:fit}.

\paragraph{Wind sea: Elfouhaily et al.\ (1997) long-wave branch (default).}
The gravity-wave part of the unified spectrum is a curvature spectrum $B_l(k)$ parameterised by the
inverse wave age $\Omega_c = U_{10}/c_p$:
\begin{align}
B_l(k) &= \tfrac12\,\alpha_p\,\frac{c_p}{c(k)}\,F_p(k),\qquad
F_p = L_{PM}\,J_p\,\exp\!\Big[-\frac{\Omega_c}{\sqrt{10}}\Big(\sqrt{k/k_p}-1\Big)\Big],\\
L_{PM} &= \exp\!\big[-\tfrac54 (k_p/k)^2\big],\qquad
J_p = \gamma^{\Gamma},\quad
\Gamma = \exp\!\Big[-\frac{(\sqrt{k/k_p}-1)^2}{2\sigma^2}\Big],\quad
\sigma = 0.08\,(1 + 4/\Omega_c^3),\\
\gamma &= \begin{cases}1.7 & \Omega_c \le 1\\ 1.7 + 6\log_{10}\Omega_c & \Omega_c>1\end{cases},
\qquad \alpha_p = 0.006\sqrt{\Omega_c},
\qquad c(k) = \sqrt{\tfrac{g}{k}\big(1+(k/k_m)^2\big)},\ k_m = 370\,\mathrm{rad\,m^{-1}} .
\end{align}
It is mapped to frequency with the deep-water dispersion relation $k = (2\pi f)^2/g$ and Jacobian
$\dd k/\dd f = 8\pi^2 f/g$: $S_{\mathrm{ws}}(f) \propto B_l(k)\,k^{-3}\,\dd k/\dd f$, with $k_p$ set from
$f_p^{(k)}$. Because the shape is renormalised to $m_0^{(k)}$, the Elfouhaily amplitude prediction is
discarded and only the shape survives: a JONSWAP-like form whose peak enhancement and width are tied
to one physical parameter, $\Omega_c\in[0.84, 5]$, with the $f^{-4}$-to-$f^{-5}$ transition of the
high-frequency slope built in. The capillary branch $B_h$ peaks at $\sim$96 Hz and contributes nothing
on the native grid. Alternatives available in the package: JONSWAP ($\gamma$), Donelan ($\gamma$,
$f^{-4}$), Pierson--Moskowitz (no parameter); these three carry a free tail exponent $n$
(Section~\ref{sec:tail}).

\paragraph{Swell: single Ochi--Hubble component (default).}
With $x = f/f_p^{(k)}$ and shape parameter $\lambda\in[1,12]$,
\begin{equation}
S_{\mathrm{sw}}(f) \propto x^{-(4\lambda+1)}\exp\!\Big[-\frac{4\lambda+1}{4}\,x^{-4}\Big] ,
\end{equation}
which reduces to Bretschneider at $\lambda=1$, sharpens as $\lambda$ grows, keeps a power-law
high-frequency tail and is strictly positive. A Gaussian in $f$ is also available; on real WW3
spectra it fails (spectral correlation $\approx0.69$ versus $0.90$), because the sum of narrow
Gaussians is a comb with gaps between partitions whereas WW3 spectra are smooth. Ochi--Hubble's
tails fill those gaps.

\subsection{Directional distribution}

Each partition uses a cosine-power distribution centred on its mean direction:
\begin{equation}
D^{(k)}(\theta_j) = \frac{\cos^{2s_k}\!\big(\tfrac12\,\Delta_j\big)}
{\sum_{j'}\cos^{2s_k}\!\big(\tfrac12\,\Delta_{j'}\big)\,\Delta\theta},
\qquad
\Delta_j = \operatorname{wrap}_{(-\pi,\pi]}\big(\theta_j - \bar\theta^{(k)}\big),
\label{eq:cos2s}
\end{equation}
normalised numerically on the 24 bins (the analytic $\Gamma$-function constant leaves percent-level
energy errors for narrow swell on so coarse a grid). The exponent follows from the archived spread by
inverting the first-moment (Kuik) relation for this family: with $\langle\cos\Delta\rangle = s/(s+1)$
and $\sigma_\theta^2 = 2(1-\langle\cos\Delta\rangle)$,
\begin{equation}
s_k = \frac{2}{\big(\sigma_\theta^{(k)}\big)^2} - 1,\qquad s_k\in[0.5, 200] .
\end{equation}
Thus each partition's first directional moment $(a_1,b_1)$ matches its archived direction and
spread in the continuum limit; on the discrete grid the match is approximate.

The wrap in \eqref{eq:cos2s} matters. Without it, a partition centred at $330^\circ$ sees the
$0^\circ$ bin as an offset of $-330^\circ$, $\cos(-165^\circ)<0$ is clipped to zero and every bin
across the $0/360$ seam receives no energy from that partition. Renormalisation hides the loss from
$\Hs$ but not from the spectral shape. This was present in the first real-data run and fixed on
2026-09-16; it moved the extratropical site from 0.869 to 0.896 spectral correlation.

\section{Fitting the shape parameters}
\label{sec:fit}

The archive provides one more integral than \eqref{eq:model} uses: $T_{m0,-1}$. Since $m_{-1}$ is
linear in the partition energies and, for a unit-$m_0$ shape, the ratio
\begin{equation}
\rho_{\mathrm{ws}}(f_p,\Omega_c) = \frac{m_{-1}}{m_0}\Big|_{S_{\mathrm{ws}}},\qquad
\rho_{\mathrm{sw}}(f_p,\lambda) = \frac{m_{-1}}{m_0}\Big|_{S_{\mathrm{sw}}}
\end{equation}
depends only on the peak frequency and the shape parameter, the constraint reads
\begin{equation}
\underbrace{\sum_k m_0^{(k)}\,\omega^{(k)}\,\rho_{\mathrm{ws}}\big(f_p^{(k)},\Omega_c\big)}_{A_{\mathrm{ws}}(\Omega_c)}
+\underbrace{\sum_k m_0^{(k)}\,\big(1-\omega^{(k)}\big)\,\rho_{\mathrm{sw}}\big(f_p^{(k)},\lambda\big)}_{A_{\mathrm{sw}}(\lambda)}
= T_{m0,-1}\,\sum_k m_0^{(k)} \equiv \mathcal{T} .
\end{equation}
The right-hand side uses the \emph{sum of the partition energies}, not the archived $m_0$, so that
the shape fit is decoupled from any energy the partitioning failed to capture (closure
$\sum_k (\Hs^{(k)})^2/\Hs^2$ is reported separately as a diagnostic; it is 1.000 for IOWAGA).
$\rho_{\mathrm{ws}}$ and $\rho_{\mathrm{sw}}$ are tabulated once on a $(f_p,\text{parameter})$ grid
(96 log-spaced $f_p$ values, 45 parameter values) on the truncated native grid, so that truncation is
included; the per-point fit is then a table lookup and a vectorised search rather than an optimiser call.

One equation for two unknowns is under-determined, so weak quadratic priors break the degeneracy:
\begin{equation}
(\Omega_c,\lambda) = \arg\min\;
\Big(\frac{A_{\mathrm{ws}}(\Omega_c)+A_{\mathrm{sw}}(\lambda)-\mathcal{T}}{\epsilon\,|\mathcal{T}|}\Big)^2
+ \Lambda\Big(\frac{\Omega_c-\Omega_c^0}{\Delta\Omega_c}\Big)^2
+ \Lambda\Big(\frac{\lambda-\lambda^0}{\Delta\lambda}\Big)^2,
\end{equation}
with $\epsilon = 0.01$ (relative tolerance on $T_{m0,-1}$), $\Lambda=0.7$, priors
$\Omega_c^0 = 1.5$, $\lambda^0 = 3$, and $\Delta(\cdot)$ the parameter range, which makes $\Lambda$
dimensionless and family-independent. The minimum is found by exhaustive search on the tabulated
grid. Where a point contains only one family the fit is exactly determined and the other prior is
inert; such points are flagged \texttt{underconstrained\_ws/sw}. Because there are exactly as many
constraints as free parameters, the priors are load-bearing, and matching $T_{m0,-1}$ does not by
itself prevent an unphysical shape---this is why the choice of swell family, not the fit, dominates
the skill.

\subsection{Tail exponent}
\label{sec:tail}

For wind-sea families with a free high-frequency exponent $n$ ($S\propto f^{-n}$; JONSWAP, Donelan,
Pierson--Moskowitz), a second stage picks $n\in\{3.5,4,\dots,6\}$ to match the archived $T_{m02}$
through $m_2 = m_0/T_{m02}^2$, which is tail-weighted and therefore carries the information the
first stage does not use. Stages alternate twice. For the Elfouhaily default the slope is intrinsic
and this stage is skipped. With either choice the reconstructed $E(f)$ runs $\sim$50\% low above
$0.3\,$Hz at all validated sites; this is the largest open error source for tail-weighted
quantities such as Stokes drift.

\section{Assembly and conserved quantities}

With $(\Omega_c,\lambda)$ fixed, the spectrum follows from \eqref{eq:model}. What is conserved:
\begin{itemize}\itemsep2pt
\item $m_0$ per partition and in total (exact, by normalisation), hence $\Hs$ up to the archive's own
partition closure;
\item $T_{m0,-1}$ to the tolerance $\epsilon$ where the priors do not dominate;
\item $T_{m02}$ approximately, when a free tail exponent is available;
\item per-partition mean direction and first-moment spread (continuum-exact, discretely approximate).
\end{itemize}
Not constrained: the whole-spectrum \texttt{dir}/\texttt{spr}, the second directional moments
$(a_2,b_2)$, and the shape of the high-frequency tail above $\sim$0.3 Hz.

\section{Validation summary}

Self-partition mode (partitions derived from the truth spectrum by a local watershed, so the numbers
are the method's intrinsic ceiling without grid-to-station sampling error), January 1993,
248 three-hourly steps per site, Elfouhaily + Ochi--Hubble:
\begin{center}
\begin{tabular}{lccc}
\toprule
site & spectral corr.\ (mean) & corr.\ p10 & $\Hs$ scatter index\\
\midrule
Tropical Eq.\ Pacific (0N 140W) & 0.917 & 0.877 & 1.1\%\\
Extratropical N Atlantic (52N 30W) & 0.896 & 0.830 & 0.15\%\\
Southern Ocean deep (58S 90W) & 0.890 & 0.838 & 0.31\%\\
N California offshore (38N 124W) & 0.901 & 0.830 & 0.45\%\\
\bottomrule
\end{tabular}
\end{center}
The spectral correlation is $\langle E_{\mathrm{rec}},E_{\mathrm{true}}\rangle/\|E_{\mathrm{rec}}\|\|E_{\mathrm{true}}\|$
with the $\Delta f\Delta\theta$ inner product. The four wind-sea families lie within 0.004 of each
other; replacing Ochi--Hubble by a Gaussian swell costs about 0.2. With the archive's own partitions read from the gridded monthly file
(\texttt{example\_reconstruct\_from\_nc.py}) the same month gives 0.885 (p10 0.821) at the
extratropical site and 0.910 (p10 0.863) at the tropical site; the gap to the self-partition numbers
is the grid-cell-to-station sampling penalty plus the difference between WW3's partitioning and the
local watershed. $\Hs$ and $T_{m0,-1}$ are reproduced to $<0.1\%$ in both modes.

\section{Known limitations}
\begin{itemize}\itemsep2pt
\item Deep water only: $k=(2\pi f)^2/g$ throughout.
\item The $a_1$-based mapping from \texttt{pspr} to $s$ has not been confirmed against the archive's
definition; if \texttt{pspr} is $a_2$-based the exponent should be obtained by numerical inversion.
\item Tail deficit above 0.3 Hz (see Section~\ref{sec:tail}).
\item \texttt{ptp} is quantised to the native frequency bins, which limits narrow swell.
\end{itemize}

\section*{References}
\begin{itemize}\itemsep1pt
\item Elfouhaily, T., B. Chapron, K. Katsaros, D. Vandemark (1997). A unified directional spectrum for long and short wind-driven waves. \emph{J.\ Geophys.\ Res.} 102(C7), 15781--15796.
\item Ochi, M. K., E. N. Hubble (1976). Six-parameter wave spectra. \emph{Proc.\ 15th Coastal Eng.\ Conf.}, 301--328.
\item Kuik, A. J., G. Ph.\ van Vledder, L. H. Holthuijsen (1988). A method for the routine analysis of pitch-and-roll buoy wave data. \emph{J.\ Phys.\ Oceanogr.} 18, 1020--1034.
\item Hanson, J. L., O. M. Phillips (2001). Automated analysis of ocean surface directional wave spectra. \emph{J.\ Atmos.\ Oceanic Technol.} 18, 277--293.
\end{itemize}
\end{document}
